\section*{Limitations}

\paragraph{Expert budget threshold.}
EC-HTE requires at least $n_{\text{expert}} / K \geq 25$ per subgroup for the hierarchical Bayesian shrinkage estimates to stabilize.
Below this threshold, per-subgroup confusion matrix estimates are too noisy for reliable inversion, and stratified MLE is preferred.
At $n_{\text{expert}} \geq 1000$, HB and stratified MLE converge to within 0.5\% on all metrics, making the hierarchical model unnecessary at large expert budgets.

\paragraph{RMSE tradeoff at small $n$.}
EC-HTE achieves lower bias than global correction but incurs higher variance when the expert budget is small relative to the number of subgroups.
In the most unfavorable configuration ($K{=}4$, extreme, $n_{\text{expert}} \leq 100$), EC-HTE CausalForest RMSE reaches 3.5$\times$ that of the naive estimator.
Practitioners should weigh whether bias or variance is more consequential for their application before adopting subgroup-stratified correction at small sample sizes.

\paragraph{Plug-in standard errors.}
The variance-adjusted CausalForest propagates outcome uncertainty through the inverse confusion matrix $\hat{C}_s^{-1}$ but does not propagate uncertainty from the confusion matrix estimation itself.
This plug-in approximation is reasonable at $n_{\text{expert}} \geq 100$, where confusion matrix estimation variance is small relative to outcome variance, but may understate total uncertainty below this threshold.
Full Bayesian propagation of confusion matrix posterior uncertainty through the causal forest splitting procedure is computationally prohibitive and left to future work.

\paragraph{Symmetric misclassification assumption.}
All experiments use symmetric error rates ($P(\hat{Z}{=}1 \mid Z{=}0) = P(\hat{Z}{=}0 \mid Z{=}1)$).
Real LLM annotators may exhibit asymmetric confusion patterns, such as higher false positive than false negative rates for certain label categories.
The bias decomposition identity (Theorem~\ref{thm:bias_identity}) generalizes to asymmetric matrices, but empirical validation under asymmetric misclassification is left to future work.

\paragraph{Subgroup structure assumed known.}
EC-HTE requires discrete, pre-defined subgroups as input.
Continuous effect modifiers or latent subgroup discovery settings require fundamentally different approaches.
The hierarchical prior provides partial robustness to noisy subgroup definitions through automatic shrinkage toward the global mean, but intersectional subgroups with many levels may exceed the $n_{\text{expert}} / K \geq 25$ budget threshold, limiting practical applicability to settings with moderate $K$.
